\section{Interpretation and related positivity methods}
\label{pos-section}

\subsection{The precise source-energy statement}

All statements in this subsection assume the hypotheses of the
Gaussian-domination theorem being used.  Let $H_0$ include the photon
terms when the full-field theorem is invoked, and let $E_\nu(q)$ be the
ground energy of $H_0+\alpha D_L(\rho,\nu)+U_{\rm nn}$ in sector $q$.
The exact completed-square statement is
\begin{equation}
 E_\nu(q)-U_{\rm nn}
 +\frac\alpha2D_L^{\rm band}(\nu,\nu)
 \geq E^{\rm vac}_{\Lambda,L}.
 \label{pos-completed-source}
\end{equation}
Here $D_L^{\rm band}(\nu,\nu)$ is the \emph{collective} source
self-energy: it contains both individual diagonals and nuclear cross
terms.  It is not just a sum of single-nucleus constants.

To state the distinction explicitly, set
\[
 G_L^{\rm band}(x)=\frac1V\sum_{0<|k|\leq2\Lambda_f}
                \frac{4\pi}{|k|^2}e^{ik\cdot x},\qquad
 U_{\rm nn}^{\rm band}=\alpha\sum_{i<j}Z_iZ_j
                  G_L^{\rm band}(R_i-R_j).
\]
Then \eqref{pos-completed-source} is equivalent to
\begin{equation}
 E_\nu(q)+\frac\alpha2\kappa_{\Lambda,L}\sum_i Z_i^2
 \geq E^{\rm vac}_{\Lambda,L}
       +U_{\rm nn}-U_{\rm nn}^{\rm band}.
 \label{pos-self-mismatch}
\end{equation}
Thus nuclei plus their individual band self-energies cannot lower the
energy below the vacuum if there is one nucleus, or if the nuclear
repulsion is also defined with the same band kernel.  For the original
point-source repulsion, the extra cross term in
\eqref{pos-self-mismatch} has no fixed sign.  Indeed,
$G_L-G_L^{\rm band}$ has zero mean, is positive near its Coulomb
singularity, and must be negative somewhere else.  Positivity of its
Fourier coefficients does not give pointwise positivity.  The
heat-kernel comparison in the many-nucleus theorem is what controls
this mismatch.

An unrestricted box-size version of the simpler physical slogan is
false for the specified periodic self-energy convention.  Choose
$y\ne0$ on the unit torus with $G_1(y)<0$, which exists because
$G_1$ has zero mean and a positive singularity.  Take two unit nuclei
at $0,Ly$ and let $L\downarrow0$ at fixed cutoffs.  Eventually only the
fermionic momentum $p=0$ remains, no nonzero photon mode remains, and
$\kappa_{\Lambda,L}=0$.  The neutral sector $q=2$ is admissible and
has electronic ground energy $2m$.  Hence
\begin{equation}
 E_\nu(2)-E^{\rm vac}_{\Lambda,L}
 =2m+\frac\alpha L G_1(y)\longrightarrow-\infty.
 \label{pos-small-box}
\end{equation}
This does not contradict the many-nucleus theorem, which assumes
$L\geq\Lambda_f^{-1}$.  It does rule out replacing
\eqref{pos-completed-source} by a stronger, unqualified statement about
the original point-nucleus Hamiltonian and individual band self-energies.

At zero source, Gaussian domination also proves
\begin{equation}
 \inf\sigma(H_0|_{Q=q})\geq
 \inf\sigma(H_0|_{Q=0})
 \quad\text{for every admissible }q.
 \label{pos-neutral-minimum}
\end{equation}
The neutral vacuum therefore minimizes the energy over all algebraic
charge sectors of the common zero-mode-deleted Hamiltonian.  Nonzero
$q$ sectors without compensating external charge are auxiliary sectors,
not isolated physical charged systems satisfying periodic Gauss law.
Neither uniqueness nor a strictly positive charge gap follows from
\eqref{pos-neutral-minimum} alone.  The stronger estimate
$E_q(0)-E_0(0)\geq m_*|q|$ under a strictly positive mass coefficient
is proved separately in Theorem~\ref{thm:mass-gap-charge}.

The compact gauge-invariant lattice model of Section~\ref{lat-section}
has a different, exact source prescription.  Its physical Gauss-law
sector retains the complete electric energy, and
Theorem~\ref{lat-main} proves $E_\nu\geq E_0$ before any nuclear
self-energy subtraction; Corollary~\ref{lat-rest-energy} strengthens
this to $E_\nu\geq E_0+m|\sum_x\nu_x|$.
Corollary~\ref{lat-linear-stability} then
subtracts the explicitly prescribed individual lattice diagonals.
For that consistently regulated model, adding those same individual
self-energies back to the subtracted energy does give an energy at
least as large as the interacting vacuum.  This does not remove the
kernel mismatch in \eqref{pos-self-mismatch} for the original
continuum point-source prescription.

\subsection{Euclidean and Gaussian positivity}

The underlying positivity mechanism has established precedents.
For a regulated Euclidean vector-coupled Dirac operator $\mathcal D$
with $\mathcal D^*=-\mathcal D$ and
$\{\gamma_5,\mathcal D\}=0$, a positive scalar mass gives
\[
 \det(\mathcal D+m)=m^{n_0}
          \prod_{\lambda>0}(m^2+\lambda^2)>0.
\]
The spectral-pairing argument is recorded explicitly in
Vafa--Witten~\cite{VafaWitten}, equation~(1).  It does not require
two identical fermion flavours.  The present proof instead works
directly with a finite-dimensional, centered Hamiltonian trace,
allows a positive spatial mass matrix, and verifies positivity after
inserting the charge-projection twist.  Its proof is supplied here
rather than inferred from a formal infinite-dimensional determinant.

The proposed Seiler reference~\cite{Seiler} requires a different
description.  Its result is an upper bound on the growth of a
lattice color-confining potential: positivity of the Hamiltonian and
Euclidean rotational symmetry prevent faster-than-linear growth.
It is not cited here as a proof of single-flavour QED determinant
positivity or of \eqref{pos-completed-source}.  A bound
$|\langle W\rangle|\leq1$ follows immediately for a bounded unitary
observable integrated against a positive probability measure, but
identifying its large-time exponent with a static-source energy also
requires a transfer-matrix statement and a fixed self-energy convention.
Those ingredients cannot be replaced by the elementary modulus bound.

Fr\"ohlich--Simon--Spencer~\cite{FSS}, Section~2, establish Gaussian
domination for periodic, nearest-neighbour classical ferromagnets on
cubic lattices.  Their shifted bond partition function is maximized
at zero shift; differentiating gives an infrared bound on the two-point
function.  Theorem~2.1 and its shifted-square formulation are the
relevant comparison, rather than a theorem for fermions or Coulomb-gauge
photons.  Miyao's finite-source result~\cite{Miyao} is closer to the
present operator setting, but uses the particular half-filled lattice
fermion structure described in Section~\ref{context:scope}.  Here the
positive centered CAR trace replaces those model-specific positivity
inputs.  Componentwise Hubbard--Stratonovich integration retains the
order of the noncommuting projected densities.

\subsection{Comparison with Majorana sign criteria}

Wei--Wu--Li--Zhang--Xiang~\cite{Wei} give two sufficient conditions for
nonnegative auxiliary-field fermion weights.  Their Theorem~1 uses
Majorana reflection-positive quadratic kernels of block form
\[
 \begin{pmatrix} A&iB\\-iB^T&A^*\end{pmatrix},
 \qquad A^T=-A,\quad B=B^*,
\]
with the stated semidefiniteness condition on $B$.  Their Theorem~2
uses a common pair of symmetry operators satisfying
$S^TVS=V^*$, $PVP^{-1}=V$, $S^2=-I$, $P^2=I$, and $PS=-SP$,
with the additional reality and symmetry conditions specified there.
These are criteria for each auxiliary-field quadratic factor; they
do not themselves identify a Coulomb Hamiltonian, its charge sector,
or its nuclear self-energy subtraction.

Li--Jiang--Yao~\cite{LiJiangYao} classify quadratic Majorana factors
by mutually anticommuting antiunitary symmetries.  The two basic
sign-free classes have symmetry-square signatures $(+,-)$ and
$(-,-)$.  This classification is within that symmetry framework;
it is not a classification of every possible nonnegative fermion
weight.  Their examples include models without particle-number
conservation.  That possibility differs from the present
electron--positron formulation, which permits pair creation while
conserving total electric charge.

The present $C_5$ identity is a pseudo-unitary spectral-pairing
condition, supplemented by a positive-mass quadratic-form condition
that excludes eigenvalues on the unit circle.  It belongs in this
broader family of positivity methods.  No unproved equivalence to a
particular Majorana criterion is required, and no priority claim for
the positivity mechanism is made.  The substantive application checked
here is its compatibility with the specified Wick correction,
canonical projection, quantized photon integration, and point-source
comparison.  None of the cited sign criteria by itself removes the
ultraviolet cutoff, repairs the gauge covariance of the sharp momentum
cutoff, or proves an extensive interacting vacuum.
